\documentclass[11pt,a4paper]{article}

\usepackage[margin=23mm]{geometry}
\usepackage[T1]{fontenc}
\usepackage[utf8]{inputenc}
\usepackage{lmodern}
\usepackage{microtype}
\usepackage{xcolor}
\usepackage{booktabs}
\usepackage{tabularx}
\usepackage{longtable}
\usepackage{array}
\usepackage{enumitem}
\usepackage{xurl}
\usepackage[hidelinks]{hyperref}
\usepackage{fancyhdr}
\usepackage{lastpage}
\usepackage{listings}
\emergencystretch=3em

\definecolor{pilotnavy}{HTML}{0B1F3A}
\definecolor{pilotblue}{HTML}{1769D2}
\definecolor{pilotteal}{HTML}{087A78}
\definecolor{pilotgreen}{HTML}{16834A}
\definecolor{pilotamber}{HTML}{B56A00}
\definecolor{pilotred}{HTML}{C52F46}
\definecolor{pilotpale}{HTML}{F3F7FC}

\newcommand{\Pilot}{\textbf{Pilot}}
\newcommand{\AION}{\textbf{AION Device Fabric}}
\newcommand{\Boardroom}{\textbf{Boardroom}}
\newcommand{\decision}[1]{%
  \par\medskip\noindent
  \colorbox{pilotblue!9}{\parbox{0.94\linewidth}{\textbf{Decided approach.} #1}}%
  \par\medskip}
\newcommand{\boundary}[1]{%
  \par\medskip\noindent
  \colorbox{pilotamber!10}{\parbox{0.94\linewidth}{\textbf{Boundary.} #1}}%
  \par\medskip}
\newcommand{\term}[1]{\textbf{#1}}

\pagestyle{fancy}
\fancyhf{}
\fancyhead[L]{\small Pilot Unified Product Architecture}
\fancyhead[R]{\small Tessaris / AION}
\fancyfoot[L]{\small Product and technical decision record}
\fancyfoot[R]{\small Page \thepage\ of \pageref{LastPage}}
\setlength{\headheight}{14pt}

\lstdefinestyle{pilotcode}{
  basicstyle=\ttfamily\footnotesize,
  backgroundcolor=\color{pilotpale},
  frame=single,
  rulecolor=\color{pilotblue!25},
  breaklines=true,
  columns=fullflexible,
  showstringspaces=false,
  keepspaces=true
}

\title{\textbf{Pilot Unified Application and Self-Hosted Intelligence Architecture}\\[5pt]
\large Bridging Personal Pilot, External Workspaces, Boardroom, Television, Mobile, and Desktop}
\author{Tessaris / AION Product and Technical Record}
\date{2 September 2026\\Decision record and implementation specification}

\begin{document}
\maketitle

\begin{abstract}
This document records the decided architecture for turning Personal Pilot,
external workspaces, Boardroom, the television experience, and the AION Device
Fabric into one coherent product.  The central proposition is one personally
owned Pilot identity, one familiar mobile application, and multiple strictly
separated spaces.  The user's intelligence and operational data run primarily
on a mother brain selected and controlled by the user or organisation.  Mobile,
television, browser, desktop, and connected devices are governed surfaces of
that intelligence rather than independent copies of it.  A person may be a
private user, household member, employee, freelancer, accountant, adviser,
director, business owner, or several of these at once.  Access is granted by
explicit role and capability, and data is never merged merely because the same
person participates in several spaces.

This is both a product decision record and a technical implementation target.
It distinguishes the already substantial Boardroom and Pilot foundations from
the unified application, networking, packaging, and production security work
that remains.  It does not claim that the native mobile application, invisible
remote tunnel, production relay, or multi-host migration system is already
complete.
\end{abstract}

\newpage
\tableofcontents
\newpage

\section{Executive Decision}

Pilot will be delivered as one product with three principal modes:

\begin{enumerate}[leftmargin=*]
  \item \term{Personal}: a private life assistant, communication centre,
        action centre, memory, service gateway, and household intelligence.
  \item \term{Workspaces}: bounded participation in an employer's, client's,
        project team's, or professional body's governed operating environment.
  \item \term{Boardroom}: the full organisational command environment for a
        person who owns, directs, administers, or is formally appointed to a
        business governance role.
\end{enumerate}

The user should experience a single, simple conversation with Pilot.  Swiping
or selecting a space changes the available context, shortcuts, permissions,
agents, records, and actions; it does not create a different application or a
different identity.

\decision{Build the Pilot Unified App next.  Do not independently complete a
second mobile product for Boardroom or continue expanding the television home
page as a separate architecture.  First establish the shared identity,
workspace, message, action, approval, receipt, and connection contracts.  Then
make mobile, television, browser, and desktop clients of those contracts.}

\subsection{Product status at this decision point}

The following percentages are planning estimates supplied during product
review, not automated completion measurements:

\begin{center}
\begin{tabularx}{\textwidth}{@{}>{\bfseries}p{36mm}p{24mm}X@{}}
\toprule
Area & Estimated maturity & Interpretation \\
\midrule
Boardroom & approximately 90\% & Broad business features and operating model exist; unified identity, mobile review, external workspace participation, and production packaging remain. \\
Personal Pilot & approximately 50\% & Core identity, tasks, messages, calendar, approvals, television, and service boundaries exist; the complete consumer experience and real connected-service qualification remain. \\
Unified mobile app & concept stage & Interactive design direction exists, but no production iOS/Android application, app-store release, background connection, or hardware-backed identity has been claimed. \\
Television surface & functional proof & A strong adoption and room-scale interface exists; it must become a client of the same person-and-space manifest rather than a separate dashboard implementation. \\
\bottomrule
\end{tabularx}
\end{center}

\section{Unified Product Model}

\subsection{One person, many spaces, different authority}

The canonical model is:

\begin{lstlisting}[style=pilotcode]
Person
  +-- Personal Pilot
  +-- Household space
  +-- Employer workspace
  +-- Client A / Marketing workspace
  +-- Client B / Accounts workspace
  +-- Advisory-board seat
  +-- Owned Boardroom (only when applicable)
\end{lstlisting}

A person does not require a business of their own to use the professional
capabilities.  They may be an employee, self-employed specialist, temporary
contractor, accountant, auditor, lawyer, adviser, supplier, finance approver,
board member, or non-executive director.  The organisation invites the person's
Pilot identity into a precisely scoped workspace and grants a signed,
revocable, expiring capability lease.

\decision{``Boardroom user'' must never be synonymous with ``business owner.''
Boardroom is the organisation's governed environment.  A workspace is the
bounded view and authority through which another person or their Pilot
participates in that environment.}

\subsection{Space types}

\begin{longtable}{@{}>{\bfseries}p{28mm}p{49mm}p{76mm}@{}}
\toprule
Space & Principal owner & Typical contents and authority \\
\midrule
\endhead
Personal & Individual & Private messages, memory, contacts, tasks, reminders, calendar, shopping, services, devices, television, files, preferences, and approvals. \\
Household & Consenting members or guardians & Shared lists, chores, home devices, deliveries, family logistics, child-safe activities, and explicitly shared memory. \\
External workspace & Organisation & Only the projects, files, agents, budgets, conversations, and approvals granted to the person's role. \\
Owned Boardroom & Organisation under the person's administration & Full business dashboards, agents, departments, evidence, decisions, configuration, roles, and governance subject to organisational policy. \\
Temporary engagement & Inviting organisation & Time-bounded access for an auditor, accountant, consultant, adviser, supplier, or project participant. \\
\bottomrule
\end{longtable}

\subsection{Role examples}

Roles are packages of capabilities, constraints, approval limits, and data
views rather than decorative titles.  Example roles include employee,
freelancer, marketing specialist, finance reviewer, accountant, auditor,
consultant, legal adviser, supplier, support operator, director, and temporary
project participant.

A freelance marketing specialist may be allowed to view an approved marketing
plan, use selected brand assets, speak with the marketing agents, submit
campaign proposals, upload deliverables, and see a bounded campaign budget.
They must not thereby receive access to payroll, banking, legal disputes,
shareholder discussion, or unrelated Boardroom deliberation.

\section{One Intelligence Across Three Primary Surfaces}

\subsection{Mobile: the private everyday control centre}

The mobile application is the primary personal identity, communication,
approval, and action surface.  It should be as understandable as a mainstream
messaging application.  It is not a compressed desktop Boardroom.

Its responsibilities are:

\begin{itemize}[leftmargin=*]
  \item conversation with Personal Pilot and authorised workspace Pilots;
  \item human and Pilot-to-Pilot messages;
  \item task, reminder, calendar, shopping, booking, and service-action cards;
  \item review, delegation, acceptance, rejection, correction, and approval;
  \item private identity, connected-service, memory, and permission control;
  \item secure control of television and other nearby AION Fabric devices;
  \item background notification and continuation when the person leaves home.
\end{itemize}

\subsection{Television: the shared room-scale surface}

The television is the adoption surface and communal display.  It is suited to
entertainment, learning, games, live companion intelligence, household status,
large-screen research, presentations, and deliberately opened work.  It is not
the default location for credentials, private messages, payment details, or
irreversible approvals.

The paired phone proves the active person.  The television then receives a
short-lived display session containing only the tiles and data that person may
show in the room.  Inactivity, explicit logout, device departure, phone
revocation, or lease expiry removes the private workspace.

\subsection{Desktop and browser: creation and detailed operation}

The desktop remains the correct surface for:

\begin{itemize}[leftmargin=*]
  \item creating a business and configuring its Boardroom;
  \item designing departments, agents, policies, workflows, and integrations;
  \item detailed financial, operational, marketing, sales, support, and people work;
  \item high-density dashboards, document production, evidence inspection, and administration;
  \item server setup, migration, recovery, update, and advanced security controls.
\end{itemize}

The browser can be a secure client of the user's mother brain, but visiting
Tessaris must not silently move the person's operational data into a
Tessaris-owned cloud account.

\decision{Mobile is for communication, situational awareness, review,
delegation, and approval.  Desktop is for setup and deep work.  Television is
for shared presence, large-screen interaction, and deliberately opened
content.}

\section{Mobile Information Architecture}

\subsection{Persistent product shell}

The mobile app maintains five destinations across Personal, Workspace, and
Boardroom modes:

\begin{enumerate}[leftmargin=*]
  \item \term{Pilot}: the main conversation, current objectives, and contextual suggestions.
  \item \term{Inbox}: human messages, Pilot-to-Pilot requests, provider messages, and workspace communications.
  \item \term{Actions}: tasks, reminders, approvals, calendar proposals, purchases, bookings, and completed receipts.
  \item \term{Spaces}: Personal, Household, employers, clients, projects, advisory seats, and owned Boardrooms.
  \item \term{Me}: identity, devices, server, accounts, memory, recovery, security, and permission management.
\end{enumerate}

The conversation composer remains in the same place in every mode.  Its active
space is always visible.  A consequential request must show which person,
organisation, account, service, and authority will be used before execution.

\subsection{Personal shortcut row}

The Personal mode may expose a compact, horizontally scrollable row containing
Calendar, Tasks, Devices, Chats, Games, Television, Learning, Shopping, Files,
and other installed capabilities.  These are entrances to contextual panels,
not permanent sections in one long controller page.

\subsection{Workspace and Boardroom shortcut rows}

An external workspace exposes only the functions granted to that role.  An
owned or administered Boardroom may expose Boardroom, Sales, Marketing,
Finance, Operations, Support, People, Products, Files, and Pilot.

The Boardroom mobile experience concentrates on:

\begin{itemize}[leftmargin=*]
  \item concise departmental dashboards;
  \item exceptions and changes since the last review;
  \item decisions requiring the person's authority;
  \item work ready to approve, reject, correct, or delegate;
  \item conversations with the Boardroom and departmental Pilots;
  \item verified receipts after an approved action;
  \item explicit escalation to the desktop for complex configuration or deep work.
\end{itemize}

\boundary{The mobile Boardroom must not attempt to reproduce every desktop
table, form, configuration panel, or business-building workflow.  Its value is
fast comprehension and governed intervention.}

\section{The Workspace Gateway}

Boardroom is integrated through a common workspace gateway rather than copied
into Personal Pilot or patched directly into the television.

\begin{lstlisting}[style=pilotcode]
Unified Pilot identity
        |
Workspace gateway
        +-- Personal Pilot
        +-- Household
        +-- External organisation workspace
        +-- Temporary professional engagement
        +-- Owned / administered Boardroom
\end{lstlisting}

The gateway resolves the active person, space, role, capabilities, data scope,
approval policy, service accounts, and presentation surface.  It returns a
signed space manifest to the client.  The client does not infer access from
which buttons happen to be visible.

\subsection{Signed workspace lease}

A workspace lease should contain at least:

\begin{itemize}[leftmargin=*]
  \item issuer organisation and target Pilot identity;
  \item workspace, role, and policy version;
  \item allowed read, propose, delegate, approve, execute, and administer capabilities;
  \item monetary, data, time, and recipient constraints;
  \item permitted service adapters and evidence classes;
  \item issue, activation, expiry, renewal, and revocation data;
  \item device or channel restrictions where appropriate;
  \item nonce, sequence, signature, and audit reference.
\end{itemize}

\subsection{Professional accountability}

Pilot may prepare work for an accountant, auditor, lawyer, finance approver, or
director, but it must not represent model output as licensed human approval.
Where professional or corporate accountability is required, the named human
must inspect and sign the exact artefact or decision.  The resulting receipt
identifies the human authority, workspace, evidence hash, action, time, and
provider result.

\section{User- and Organisation-Owned Intelligence}

\subsection{Mother brain principle}

The mother brain is the authoritative runtime for cognition, memory, policy,
identity, service credentials, action coordination, and audit.  It may run on:

\begin{itemize}[leftmargin=*]
  \item a supported Mac, Windows computer, or Linux computer;
  \item a NAS or home server;
  \item a dedicated Pilot appliance;
  \item the user's own virtual private server;
  \item a business-controlled private cloud or on-premises server;
  \item a selected third-party provider under the user's account;
  \item a hybrid design with a private primary and an encrypted availability node.
\end{itemize}

\decision{Tessaris distributes and updates the software but does not become the
default owner of the customer's memory, documents, credentials, operational
database, or Boardroom evidence.}

\subsection{Deployment profiles}

\begin{longtable}{@{}>{\bfseries}p{27mm}p{40mm}p{45mm}p{36mm}@{}}
\toprule
Profile & Best suited to & Advantages & Principal trade-offs \\
\midrule
Home local & Families and privacy-conscious individuals & Local data, local device access, low recurring infrastructure cost & Availability depends on home power, network, and maintenance \\
Dedicated appliance & Non-technical households & Predictable installation, hardware, recovery, and support & Hardware purchase and product-support obligation \\
Personal VPS & Technical individuals and mobile users & Continuous availability and user-controlled server account & Remote device access may require a home gateway; user pays provider \\
Business on-premises & Cautious or regulated organisations & Direct infrastructure control, internal identity integration, local evidence & Requires professional administration and support \\
Business private cloud & Distributed organisations & Availability, scaling, managed infrastructure options & Cloud cost, configuration, jurisdiction, and provider dependency \\
Hybrid & Households or businesses needing local control and continuity & Local-first device access plus remote availability & Replication, key management, and conflict recovery are more complex \\
\bottomrule
\end{longtable}

\subsection{Simple and advanced setup}

Most users should see only:

\begin{center}
\textbf{Install Pilot} $\rightarrow$ \textbf{Connect my phone} $\rightarrow$
\textbf{Choose what Pilot may access} $\rightarrow$ \textbf{Start}
\end{center}

Advanced users and organisations may inspect or configure host location,
encryption keys, backup target, retention, model packs, relay policy, network
routes, log export, update channel, service scopes, and disaster recovery.
These controls must exist without becoming mandatory onboarding screens for an
ordinary user.

\section{Secure Connection Architecture}

\subsection{Product language}

The user action is \emph{Connect to my Pilot}, not \emph{configure a VPN}.
An encrypted private network may implement the connection, but its complexity
belongs below the product surface.

\subsection{Pairing flow}

\begin{enumerate}[leftmargin=*]
  \item The mother brain creates a short-lived, single-use pairing challenge.
  \item A trusted local surface displays a QR code and human-verifiable code.
  \item The phone creates a device key in platform-protected storage.
  \item The phone and mother authenticate the challenge and establish each
        other's public identity.
  \item The mother issues a revocable device certificate and bounded initial lease.
  \item The phone confirms the server fingerprint and recovery options.
  \item Local and remote connection paths are tested without exposing service secrets.
  \item The event is recorded in the owner-visible device and audit registers.
\end{enumerate}

\subsection{Connection selection}

\begin{lstlisting}[style=pilotcode]
1. Direct local connection over the trusted home/business network
2. Direct encrypted peer-to-peer connection across the Internet
3. End-to-end encrypted relay when NAT/firewall traversal fails
4. Honest offline state with queued local work where safe
\end{lstlisting}

The mother brain remains the trust anchor.  A relay forwards opaque encrypted
traffic and must not possess the content-encryption keys.  Connection discovery
may be operated by Tessaris, self-hosted, or supplied by another provider.

\subsection{Role of tessaris.ai}

The public site may provide software distribution, documentation, account and
licence management, optional connection rendezvous, update metadata, and a
secure web client bootstrap.  The intended web sequence is:

\begin{enumerate}[leftmargin=*]
  \item visit Tessaris and choose \emph{Connect to my Pilot};
  \item authenticate the personal device;
  \item select a registered mother brain;
  \item negotiate an encrypted session with that mother brain;
  \item load the authorised interface and data from the user's environment.
\end{enumerate}

The web domain must not be treated as permission to centralise customer data.
The browser client must verify the target mother identity and active space.

\subsection{Implementation progression}

The first production-capable release may use a proven WireGuard- or
Tailscale-style private network while the native Pilot application owns the
identity and user experience.  A later embedded tunnel may remove the separate
network application and account.  Any replacement must match or exceed the
tested security, roaming, key rotation, revocation, and recovery properties of
the initial system.

\section{Identity, Authority, and Privacy Boundaries}

\subsection{Identity layers}

Pilot separates:

\begin{itemize}[leftmargin=*]
  \item person identity;
  \item device possession;
  \item biometric or operating-system approval;
  \item mother-brain identity;
  \item household membership;
  \item organisation membership;
  \item role and capability lease;
  \item active presentation session;
  \item service-account consent.
\end{itemize}

Possession of one layer does not imply all others.  For example, control of a
television remote does not establish a private identity, and membership in a
marketing workspace does not grant access to the whole Boardroom.

\subsection{Data ownership}

The individual owns Personal Pilot data.  The organisation owns its Boardroom
and workspace records.  A shared household owns only those objects explicitly
created or moved into the household scope.  Service providers retain whatever
records their own terms and systems require; Pilot stores references and
normalised receipts without copying secrets unnecessarily.

\subsection{Cross-space transfer}

Moving information between spaces is an explicit, inspectable operation:

\begin{quote}
``Send this personal travel receipt to my company expenses.''
\end{quote}

Pilot displays the source, destination, exact payload, retention consequence,
and intended recipient before transferring it.  It records a receipt in both
spaces without manufacturing a hidden shared memory.

\boundary{The Boardroom cannot silently read personal messages, calendars,
household history, other clients, private files, or personal memory.  Personal
Pilot cannot copy organisational evidence or expose it on a television merely
because the same person can access both spaces.}

\subsection{Availability without calendar disclosure}

Where a workspace needs scheduling information, Personal Pilot may answer a
bounded question such as \emph{available} or \emph{unavailable} for a proposed
period.  It does not disclose private event names, participants, locations, or
notes unless the person explicitly authorises that disclosure.

\section{Unified Communication and Action Protocol}

\subsection{Canonical flow}

\begin{lstlisting}[style=pilotcode]
Request made
  -> person and active space resolved
  -> intent classified and policy checked
  -> exact action proposal prepared
  -> missing private details collected privately
  -> approval requested when required
  -> recipient or provider receives the action
  -> result independently verified
  -> receipt returned to the sender and correct space
\end{lstlisting}

\subsection{Message and action types}

The unified stream carries typed objects rather than undifferentiated text:

\begin{itemize}[leftmargin=*]
  \item human message;
  \item Pilot-to-Pilot message;
  \item task proposal, acceptance, decline, snooze, and completion;
  \item reminder and location-aware reminder proposal;
  \item calendar create, reschedule, cancel, and availability proposal;
  \item document review and comment request;
  \item Boardroom decision or professional sign-off request;
  \item email, messaging, or call draft;
  \item purchase, booking, or service-action proposal;
  \item device-control request;
  \item emergency or priority communication;
  \item progress, failure, correction, and verified completion receipt.
\end{itemize}

Every object carries a sender, intended recipient, originating space, allowed
destination, privacy label, creation time, expiry where appropriate, state,
policy decision, and audit reference.

\subsection{Pilot recipient and non-Pilot recipient}

If the recipient has Pilot, the request arrives as a structured object in the
recipient's own mother brain and inbox.  The recipient's Pilot can explain,
schedule, negotiate, decline, or request clarification within its authority.

If the recipient does not have Pilot, the sender may approve a delivery through
an officially authorised email, SMS, WhatsApp-equivalent, share link, or other
adapter.  The system must distinguish \texttt{prepared}, \texttt{approved},
\texttt{submitted}, \texttt{delivered}, \texttt{accepted}, and
\texttt{completed}.  It must not claim a message or task was received merely
because a model generated it.

\subsection{Representative journeys}

\begin{itemize}[leftmargin=*]
  \item ``Send Becca a task to collect the dry cleaning.''
  \item ``Email James and ask for the revised proposal.''
  \item ``Tell my marketing freelancer to review this campaign.''
  \item ``If they have not replied tomorrow, remind me.''
  \item ``Move this message into the Acme marketing workspace.''
  \item ``Show the proposal on the television, but keep the financial appendix private.''
  \item ``Ask Finance Pilot for the strongest objection before I approve this.''
\end{itemize}

\section{Boardroom and External Workspace Integration}

\subsection{Boardroom as a capability provider}

Boardroom exposes agents, departments, dashboards, projects, files, evidence,
decisions, approvals, and messages through the workspace gateway.  The unified
app consumes these resources according to the active workspace lease.  It does
not duplicate the Boardroom database or reimplement the organisation's policy.

\subsection{Mobile Boardroom behaviour}

The mobile Boardroom opens with:

\begin{itemize}[leftmargin=*]
  \item a visible organisation and role;
  \item compact department navigation;
  \item a concise status briefing;
  \item a small number of meaningful operational indicators;
  \item exceptions, decisions, approvals, and delegated work;
  \item a persistent Boardroom Pilot conversation composer.
\end{itemize}

Selecting Sales, Marketing, Finance, Operations, Support, People, Products, or
Files changes the evidence and departmental Pilot context.  It does not grant
new authority.  Complex setup or detailed work offers an explicit
\emph{Continue on desktop} action.

\subsection{External participant flow}

\begin{enumerate}[leftmargin=*]
  \item An authorised organisation member invites a person's Pilot identity.
  \item The invitation declares the organisation, role, requested capabilities,
        data classes, duration, and approving authority.
  \item The individual reviews and accepts the invitation privately.
  \item Both mother brains record the relationship and signed workspace lease.
  \item The workspace appears in the person's Spaces view.
  \item Work is requested, delivered, reviewed, and receipted through the
        organisation's policy boundary.
  \item Expiry, role change, project completion, or revocation removes access
        without deleting the person's unrelated Pilot identity.
\end{enumerate}

\section{Television Integration}

The television home page is rebuilt as a client of the same signed surface
manifest.  The manifest may expose Personal Pilot, Tasks, Calendar, Messages,
Files, Games, AI TV, Learning, Shopping, Home/IoT, Workspace, and Boardroom.
Only capabilities authorised for the active person and safe for a shared screen
are rendered.

\subsection{Taking control}

\begin{enumerate}[leftmargin=*]
  \item A trusted phone discovers or selects the television.
  \item The phone proves device possession and active person.
  \item The mother checks household/device trust and issues a short-lived
        television presentation lease.
  \item The television displays the person's name, active space, expiry, and
        permitted tiles.
  \item Consequential approval remains on the private phone.
  \item Explicit logout, inactivity, departure, phone revocation, or lease
        expiry clears the private screen state.
\end{enumerate}

\subsection{Shared-screen minimisation}

The television normally receives counts, summaries, redacted titles, or a
deliberately opened artefact.  It does not receive reusable private keys,
service tokens, complete private inboxes, payment credentials, or unrestricted
Boardroom memory.

\section{Core Domain Contracts}

The following provider-independent objects should be versioned before large
client implementation begins:

\begin{longtable}{@{}>{\bfseries}p{33mm}p{120mm}@{}}
\toprule
Object & Minimum purpose \\
\midrule
Person identity & Stable human-controlled identifier, display claims, guardian relationship where applicable, recovery and revocation state. \\
Device identity & Public key, platform, possession proof, trust state, capabilities, last-seen evidence, and revocation. \\
Space & Personal, household, workspace, temporary engagement, or Boardroom ownership and policy anchor. \\
Membership & Person-to-space relationship, role, start, expiry, invitation, acceptance, and revocation. \\
Capability lease & Signed authority, constraints, sequence, expiry, policy version, issuer, subject, and permitted surfaces. \\
Conversation & Participants, owning space, message sequence, privacy, retention, and continuation state. \\
Action proposal & Exact requested change, risk, evidence, missing fields, approval requirements, and expiry. \\
Approval & Person, trusted device, exact proposal hash, biometric/OS assertion where available, time, and result. \\
Receipt & Submission, provider response, independent verification, correction, and final state without provider secrets. \\
Surface manifest & Active person, device, space, permitted modules, redaction rules, expiry, and signature. \\
Memory item & Owner, scope, provenance, confidence, retention, correction, export, and deletion controls. \\
Service connection & Persona or organisation owner, provider, scopes, credential reference, consent, health, and revocation. \\
\bottomrule
\end{longtable}

\subsection{Minimum state language}

Clients and agents should share precise states such as
\texttt{draft}, \texttt{needs\_details}, \texttt{awaiting\_approval},
\texttt{approved}, \texttt{submitted}, \texttt{delivered},
\texttt{awaiting\_recipient\_acceptance}, \texttt{accepted},
\texttt{declined}, \texttt{completed}, \texttt{failed},
\texttt{cancelled}, and \texttt{expired}.  Interfaces may translate these into
plain language but must not collapse them into a misleading success state.

\section{Security and Threat Model}

\subsection{Principal threats}

The design must explicitly test:

\begin{itemize}[leftmargin=*]
  \item stolen or replaced phone;
  \item malicious QR code or mother-brain impersonation;
  \item replayed invitation, lease, activation, or approval;
  \item compromised relay or connection-discovery service;
  \item cross-space data leakage and confused-deputy actions;
  \item malicious workspace administrator requesting excessive access;
  \item prompt injection from messages, files, web pages, or Boardroom evidence;
  \item service-token theft and excessive OAuth scope;
  \item household member using an unattended television session;
  \item lost server, failed update, corrupt database, or unavailable home network;
  \item a model presenting an unverified action as completed;
  \item abusive contact, child-safety, harassment, and emergency misuse.
\end{itemize}

\subsection{Required controls}

\begin{itemize}[leftmargin=*]
  \item platform-protected device keys and signed, short-lived challenges;
  \item certificate pinning or equivalent verified mother-brain identity;
  \item least-privilege, expiring, revocable capability leases;
  \item independent Personal, Household, Workspace, and Boardroom encryption and policy domains;
  \item explicit private approval for consequential actions;
  \item separate service credentials by person or organisation;
  \item end-to-end encryption across any relay;
  \item redacted audit records and inspectable action receipts;
  \item safe cancellation, timeout, replay protection, and idempotency;
  \item encrypted backup, tested recovery, key rotation, and lost-device revocation;
  \item no model access to raw keys, passwords, card numbers, or unrestricted tokens.
\end{itemize}

\section{Operations, Updates, Migration, and Recovery}

Self-hosting is credible only if ordinary people can operate it.  The runtime
must provide:

\begin{itemize}[leftmargin=*]
  \item guided installation with no terminal requirement;
  \item health monitoring and non-technical explanations;
  \item signed updates, staged activation, health verification, and rollback;
  \item automatic restart and bounded recovery after power or network loss;
  \item encrypted backup with owner-visible status;
  \item export, selective deletion, and complete mother migration;
  \item lost-phone, lost-server, and compromised-key recovery;
  \item a redacted support bundle that excludes memories, credentials, and private message bodies;
  \item a compatibility report covering host, network, devices, and provider adapters.
\end{itemize}

Migration must transfer encrypted user-owned state to a new authorised mother,
rotate the active endpoint and relevant keys, verify the new instance, and
retire the old authority.  It must not silently create two writable primaries.

\section{Adoption and Network Effects}

The television remains the visible acquisition wedge.  The unified mobile app
creates daily utility, and workspace invitations create professional adoption.

\begin{enumerate}[leftmargin=*]
  \item A household installs Pilot for the television, games, learning, or live intelligence.
  \item The paired phone becomes useful for tasks, messages, reminders, approvals, and continuation.
  \item A person sends a useful task, result, or request to another person.
  \item A business invites an employee, freelancer, accountant, or adviser into a workspace.
  \item The invited person installs Pilot to collaborate without surrendering their personal identity.
  \item They discover that the same privately owned Pilot can manage their personal life and other workspaces.
  \item They may later create a household, workspace, or Boardroom of their own.
\end{enumerate}

\decision{The viral loop is utility, not advertising: a real task, decision,
message, Moment, family request, or professional invitation gives another
person a reason to connect.}

\section{Implementation Programme}

\subsection{Phase A: canonical contracts}

Version the Person, Device, Space, Membership, Capability Lease, Conversation,
Action, Approval, Receipt, Surface Manifest, Memory, and Service Connection
schemas.  Implement policy tests before UI integration.

\subsection{Phase B: unified app shell}

Build the Personal/Workspace/Boardroom switcher; Pilot, Inbox, Actions, Spaces,
and Me navigation; contextual shortcut rows; conversation composer; action
cards; and clear active-person, active-space, and connection-state indicators.

\subsection{Phase C: secure pairing and remote access}

Implement QR pairing, device keys, certificates, local discovery, direct remote
connection, encrypted relay fallback, key rotation, revocation, and recovery.

\subsection{Phase D: unified communication loop}

Implement structured Pilot-to-Pilot messaging, recipient acceptance, task and
calendar conversion, provider fallbacks, follow-up reminders, delivery states,
and receipts.

\subsection{Phase E: Personal Pilot integration}

Connect personal identity, memory, tasks, calendar, contacts, service accounts,
household, television, devices, saved items, shopping, booking, and files.

\subsection{Phase F: Boardroom and workspace integration}

Expose Boardroom functions through the workspace gateway.  Implement invitation,
role, lease, departmental briefing, review, delegation, approval, professional
sign-off, and desktop continuation.

\subsection{Phase G: television convergence}

Replace hard-coded television tiles with the signed surface manifest.  Verify
private takeover, timeout, logout, redaction, personal and workspace views, and
phone-only approval.

\subsection{Phase H: production clients and distribution}

Package signed native iOS and Android clients, secure browser access, desktop
clients, server installers, update and rollback, diagnostics, backup, and
migration.  Complete clean-machine and ordinary-user field qualification.

\section{Acceptance Programme}

The unified bridge is not complete until all of the following journeys pass
with independent evidence:

\begin{enumerate}[leftmargin=*]
  \item A new household installs a mother brain and pairs a phone without using a terminal.
  \item The phone works locally, leaves Wi-Fi, reconnects remotely, and returns without losing identity or duplicating an action.
  \item A user switches between Personal, external Workspace, and owned Boardroom without cross-space leakage.
  \item Personal Pilot sends a structured task to another Pilot; the recipient accepts; both see consistent state and receipts.
  \item A non-Pilot recipient receives an approved provider message, and the sender sees an honest delivery state.
  \item An organisation invites a freelancer with marketing-only authority; forbidden Finance and People data remain inaccessible.
  \item An accountant receives a bounded review package, signs an exact result, and loses access when the engagement expires.
  \item A Boardroom owner reviews a dashboard, delegates a task, approves a decision privately, and receives provider verification.
  \item A phone activates a television workspace; private content stays redacted; inactivity and explicit logout clear the session.
  \item A lost phone is revoked and cannot reconnect, approve, activate a television, or renew a workspace lease.
  \item A mother brain is backed up, migrated, verified, and the old writable authority is retired.
  \item A failed update rolls back without losing identity, messages, tasks, permissions, or audit continuity.
\end{enumerate}

\section{Non-Goals and Honest Boundaries}

This phase does not require:

\begin{itemize}[leftmargin=*]
  \item reproducing full desktop Boardroom administration on mobile;
  \item forcing every user to understand VPNs, certificates, ports, or server administration;
  \item moving customer memory and business evidence into a Tessaris-owned central cloud;
  \item allowing an LLM to invent permissions, select a private identity, or execute outside a signed lease;
  \item claiming biometric identity until a signed native client and hardware-backed proof are field-qualified;
  \item claiming message delivery, professional approval, purchase, booking, or device action without a verified external receipt;
  \item merging a person's Personal Pilot memory with employer or client data;
  \item treating a visible mobile button as proof that the underlying capability exists.
\end{itemize}

\section{Final Product Definition}

\begin{center}
\large\textbf{One personally owned Pilot identity that can securely enter many
homes, teams, projects, and Boardrooms without merging their data or confusing
their authority.}
\end{center}

The television gets Pilot into the home and demonstrates immediate value.  The
mobile application turns Pilot into an everyday communication and action
system.  Personal Pilot builds trust and continuity.  External workspaces let
employees and specialists participate safely in other organisations.
Boardroom provides the complete governed business environment.  Desktop and
browser clients support detailed creation and operation.  The user- or
organisation-selected mother brain preserves ownership of intelligence and
data across all of them.

\appendix
\section{Initial Product Screen Contract}

\subsection{Personal mode}

Header: Pilot identity and active Personal space.  Shortcut row: Calendar,
Tasks, Devices, Chats, Games, Television, Learning, Shopping, and Files.
Stream: conversations, action cards, household requests, provider results, and
receipts.  Composer: \emph{Ask Pilot anything}.  Bottom navigation: Pilot,
Inbox, Actions, Spaces, and Me.

\subsection{External workspace mode}

Header: organisation, workspace, and role.  Shortcut row: only granted
functions.  Stream: work requests, messages, reviews, deadlines, delegated
actions, and receipts.  Composer: active workspace Pilot.  Personal messages
and other client data are absent.

\subsection{Boardroom mode}

Header: organisation, active Boardroom role, and connection state.  Shortcut
row: Boardroom, Sales, Marketing, Finance, Operations, Support, People,
Products, Files, and Pilot as permitted.  Stream: briefing, exceptions,
decisions, approvals, delegations, and receipts.  Detailed configuration offers
desktop continuation.

\section{Decision Register}

\begin{longtable}{@{}>{\bfseries}p{32mm}p{121mm}@{}}
\toprule
Decision & Recorded outcome \\
\midrule
Product structure & One Pilot product with Personal, Workspace, and Boardroom modes. \\
Primary mobile model & Messaging-simple conversation plus structured action cards, not a miniature desktop. \\
Boardroom mobile role & Review, briefing, delegation, approval, and intervention; setup and deep work remain desktop-first. \\
Identity & One person may participate in many spaces with different authority. \\
External work & Employees, freelancers, advisers, accountants, auditors, directors, and suppliers use bounded workspaces. \\
Hosting & User- or organisation-selected mother brain; local, appliance, private server, private cloud, or hybrid. \\
Tessaris cloud role & Distribution, optional rendezvous/relay, updates, and web bootstrap; not default ownership of customer intelligence or data. \\
Remote connection & Local direct first, encrypted peer-to-peer second, end-to-end encrypted relay fallback. \\
Television & Shared room-scale client using short-lived person-and-space manifests. \\
Data transfer & Explicit, inspectable, receipted cross-space action; no hidden memory merging. \\
Execution & Model output proposes; policy, private authority, provider adapter, verification, and receipt determine completion. \\
Next programme & Pilot Unified App -- Identity, Spaces, Inbox, Actions, secure connection, then Personal/Boardroom/TV integration. \\
\bottomrule
\end{longtable}

\end{document}
